\section{Discussion and Conclusion}

\noindent\textbf{What simulation buys, and for which task.}
\emph{Pose} needs coverage of the object's contact silhouettes over the pose space, which the geometry stage provides regardless of appearance: a physical renderer fitted to one no-contact frame makes the pose of a new object estimable, covering the rotation range matters more than rendering quality, and simulation keeps pose stable across sessions. \emph{Depth and normals} depend on the appearance inside the contact, which only a learned renderer reproduces, and only after $\approx\!25$--$50$ labeled contacts per object; below that nothing leakage-free helps, above $\approx\!100$ real data alone is already good. Choose by task: for pose of known meshes, calibrate a physical renderer from a background frame and cover the rotations; for dense geometry, budget the real pairs the learned renderer needs, condition it on the session background so that new sessions and objects cost one frame, and account for those pairs when reporting.

\noindent\textbf{Reporting.}
Learned-renderer results should state the renderer's real training set and keep it disjoint from the downstream test frames. On our data the leaky protocol overstates the depth benefit by $2$--$4\times$ where the pipeline is most attractive, and manufactures a $30\%$ gain on unseen objects that no leakage-free renderer reproduces. The contact-region fidelity proxy makes the strict protocol cheap: nineteen renderer variants were ranked without downstream training, and the ranking held.

\noindent\textbf{Limitations.}
Sensor-A data have one session per object, so object identity and session appearance are confounded in the leave-object-out study; the second session and sensor separate them only partially (re-estimated poses, a different capture protocol). Depth references are mesh projections, not membrane measurements. The downstream model is one frozen-encoder architecture; a trainable encoder may extract more from synthetic data or be more sensitive to its artifacts. The physical renderer uses a geometric gel model; a mechanical model could improve the contact shading we identify as the bottleneck.

\noindent\textbf{Conclusion.}
We compared a physically calibrated and a learned tactile renderer on two fingertip sensors, three tasks and four transfer protocols, under a rule that charges each renderer for the real data it consumes and forbids it from seeing the test frames. Simulation transfers tactile pose estimation to objects without real data and stabilizes it across sessions with a renderer calibrated from a single frame; dense depth and normals improve only with a learned renderer and a moderate real budget, on both sensors, and the benefit is overstated by up to $4\times$ when the renderer sees the test distribution. A background-conditioned learned renderer and rotation coverage extend the useful range at no extra labeled data, while physics-guided, illumination-normalized and difference-image renderers do not; the few-shot gap lies in the shading inside the contact.
